\subsubsection{Actividades [Pendientes: 1. seleccionar sector, 2. confirmar datos y acceso HPC, 3. verificar revisión ejecutable de FDS]}\label{sec:inicio-actividades}
Se consolidarán las necesidades de las entrevistas y se delimitarán requisitos y exclusiones con el equipo y el director, relacionándolos con los objetivos de la sección~\ref{sec:objetivos}. A continuación, se seleccionará el sector a partir de su pertinencia y de la disponibilidad defendible de terreno, coberturas y meteorología, y se preparará un inventario que distinga fuentes identificadas, archivos disponibles, datos transformados y entradas efectivamente comprobadas. Para cada parámetro se conservarán procedencia, unidad y condición de aplicación, evitando completar las brechas con valores de ejemplo. Al mismo tiempo, se describirán mediante los modelos necesarios la organización de componentes, los contratos de intercambio y los metadatos que conectarán la preparación científica con la consulta, dentro de los entregables E1, E2 y E3.

Se realizarán prácticas de FDS que permitan verificar la revisión utilizada, su construcción o instalación y la ejecución de un caso controlado, antes de asumir que el repositorio identificado proporciona una configuración lista para el territorio. También se comprobarán las condiciones de acceso al recurso HPC previsto, ZINE, y las modalidades de ejecución que realmente permita la combinación de simulador e infraestructura. La oferta de recursos y capacitación conocida por el equipo será tratada como antecedente de viabilidad y no como una asignación ya concedida. Finalmente, se prepararán los primeros criterios y casos de prueba del flujo, incluyendo la identificación de condiciones no representadas y la revisión de coherencia espacial, para que las comprobaciones acompañen el desarrollo desde esta fase.
